\documentclass[11pt, ngermanm, titlepage]{article}
\usepackage[ngerman]{babel}
\usepackage{fancyhdr}
\usepackage{multicol}
\usepackage{a4wide}
\usepackage{tabu, tabularx}
\usepackage{graphicx}
\usepackage{fontspec}

\setmainfont{Arial}
\newfontfamily\jubilo[Ligatures=TeX]{Monotype Corsiva}

\setlength{\columnsep}{1cm}
\setlength{\columnwidth}{10cm}
\setlength{\extrarowheight}{15pt}

\begin{document}
	\pagenumbering{gobble}
	\begin{titlepage}
	\centering
	\includegraphics[width=0.8\textwidth]{img/cvlogo.png}\par\vspace{2cm}
	
	{\scshape\Large Festliches Sommerkonzert\par}
	\vspace{1.5cm}
	{\fontsize{50}{60}\bfseries\jubilo In dulci jubilo\par}
	\vspace{1.5cm}
	{\scshape\LARGE Camerata Vocale München\par}
	\vspace{1cm}
	{\Large\itshape Leitung: Clayton Bowman\par}
	\vfill
	21.07.2019\par
	Erlöserkriche München
	
	\vfill
	
	% Bottom of the page
	{\scriptsize Mitglied im\par}
	\includegraphics[width=0.25\textwidth]{img/vdkc_logo_klein.jpg}\par\vspace{2cm}


	\end{titlepage}

	\pagebreak

	
	\section*{Programm}
	\begin{tabularx} \textwidth {lX}
		Claude Debussy (1862 - 1918) & Trois Chansons: \newline I. Dieu! qu'il la fait bon regarder! \newline II. Quant j'ai ouy le tabourin \newline III. Yver, vous n'estes qu'un villain;
		\\
		Edward Elgar (1857 - 1934) & Two Choral Songs, Op. 71: \newline 1. The Shower \newline 2. The Fountain
		\\
 		 & Death on the Hills, Op. 72
 		\\
 		 & Two Choral Songs, Op. 73: \newline 1. Love's Tempest \newline 2. Serenade
		\\
		Igor Stravinsky (1882 - 1971) & Mass: \newline Kyrie \newline Gloria \newline Credo \newline Sanctus \newline Agnus Dei
	\end{tabularx}
	\pagebreak

	\begin{multicols}{2}
	\section*{Andreas Hammerschmidt}
	\paragraph{(geboren 1611 oder 1612 in Brüx, Böhmen; gestorben 1675 in Zittau)\newline}
	Der deutsche Organist und Komponist Andreas Hammerschmidt (auch „-schmied“ oder „-schmiedt“) floh mit seiner Familie vor der forcierten Rekatholisierung in Böhmen und gelangte so nach Freiberg im Kurfürstentum Sachsen, wo er vermutlich seine musikalische Ausbildung erhielt. Nach zwei vorausgehenden Organistenstellen ließ er sich schließlich 1639 als Organist in der damals reichen Stadt Zittau nieder, wirkte dort bis zu seinem Lebensende und gelangte zu gesellschaftlichem Ansehen und überdurchschnittlichem materiellem Wohlstand.
	
	Als Komponist des Frühen bis Mittleren Barock ist er in die Gruppe der evangelisch-lutherischen Kirchenkomponisten wie Heinrich Schütz und Johann Sebastian Bach einzuordnen. Sein kompositorisches Schaffen umfasst unter anderem Lieder, Kantaten, Motetten und Instrumental- wie auch Vokalkompositionen. Im Jahre 1757 vernichtete der große Stadtbrand in Zittau leider einen Großteil der Quellen über Hammerschmidt.
	
	\textit{Quelle: Deutsch- und englischsprachige Wikipedia}

	
	\section*{Thomas Weelkes}
	\paragraph{(geboren 1576 in Elsted, West Sussex; gestorben 1623 in London)\newline}
	Der englische Komponist Thomas Weelkes wurde 1598 zum Organisten des Winchester College ernannt und wechselte 1601 in die Kathedrale von Chichester. Dort wurde er jedoch später entlassen, da er an der Orgel getrunken hatte und während des Gottesdienstes schmutzige Worte benutzte. Weelkes, ein Meister der Wortmalerei, war als Madrigalist bekannt, dessen zweiter Band (1600 veröffentlicht) eine der wichtigsten Sammlungen in der englischen Madrigal-Tradition ist.

	\textit{Quellen: emmanuelmusic.com,	Wikipedia}
	
	
	\section*{Francis Poulenc}
	\paragraph{(geboren 1899, Paris; gestorben 1963, ebenda)\newline}
	Francis Jean Marcel Poulenc wuchs in einem musikalischen Haushalt auf und erhielt ersten Klavierunterricht durch seine Mutter. Nach weiterer Ausbildung als Pianist blieb das Klavier zunächst zentral und dominiert seine frühen Werke. Nicht nur Stravinsky, Chevalier und das französische Vaudeville beeinflussten ihn, sondern auch der Freund und Mentor Erik Satie, der sich dem damaligen französischen musikalischen Mainstream entzog, sowie seine Mitgliedschaft im Komponistenzirkel „Les Six“, der den Impressionismus zugunsten einer größeren Einfachheit und Klarheit ablehnte. So stand Poulenc beispielsweise Ravels musikalischen Ansichten kritisch gegenüber, wenngleich er ihn als Person sehr respektierte. Poulencs Bewusstsein für seine Insuffizienz in Bezug auf das technische Niveau seiner weitestgehend autodidaktisch erworbenen kompositorischen Fertigkeiten bewog ihn später zu Unterricht bei Charles Koechlin. Nach ab den 1920er Jahren – auch als Pianist – eintretendem internationalem Erfolg führten 1936 der Unfalltod eines Freundes und der Besuch des Wallfahrtsortes Rocamadour dazu, dass sich der Komponist intensiv auf den katholischen Glauben zurückbesann, mit der Komposition geistlicher Werke begann und eine stilistisch ernstere Komponente entwickelte. Die letzterer entspringenden Werke erhielten erst in den 1950er Jahren größere Anerkennung.
	
	Im Privatleben stellten der (während der Nazibesetzung auch politische) Umgang mit seiner Homosexualität sowie wiederkehrende, seine Schaffenskraft beeinflussende Depressionen Herausforderungen dar.
	
	Poulenc selbst verstand seinen Schaffensschwerpunkt in der Komposition von Opern. Kritiker sehen insbesondere die Dualität seiner stilistischen Pole: leichtherzig bis pietätlos versus ernsthaft-gewichtig. Das Gesamtwerk gilt als weitestgehend diatonisch und melodiezentriert und umfasst Bühnenwerke, Filmmusik, Geistliche Werke, weltliche Chorwerke, Kammermusik, zahlreiche Lieder sowie Klavier- und Orchesterwerke.
	
	\textit{Quelle: Deutsch- und englischsprachige Wikipedia}
	
	\section*{Michael Praetorius}
	\paragraph{(geboren 1571 in Creuzburg bei Eisenach; gestorben 1621 in Wolfenbüttel)\newline}
	Michael Praetorius, eigentlich Michael Schulteis, war ein deutscher Komponist, Organist, Hofkapellmeister und Gelehrter im Übergang von der Renaissance- zur Barockzeit. Sein Vater war Pfarrer, der in Wittenberg bei Martin Luther studiert hatte. Auch die beiden älteren Brüder wurden Pfarrer, und so schrieb sich auch Michael Praetorius schon mit 11 Jahren an der Universität in Frankfurt/Oder für Theologie ein. Als der Vater und die Brüder starben, war niemand mehr da, der den 16-jährigen finanziell unterstützen konnte. Er wechselte auf die Orgelbank, um seinen Lebensunterhalt zu verdienen. Auch ohne geregelte musikalische Ausbildung erwarb er sich schnell einen überregionalen Ruf. So kam er an den Hof des musikliebenden Herzogs Heinrich Julius zu Braunschweig. 
	
	Am Hof in Wolfenbüttel und Gröningen stieg er zum Hofkomponisten auf und entfaltete in kurzer Zeit eine ungeheuer fruchtbare Nebentätigkeit als Komponist und Herausgeber. Er war ein glühender Anhänger der aufregenden musikalischen Neuerungen aus Italien am Übergang von der Renaissance zum Barock. Leider konnte er nie selbst nach Italien reisen, aber durch Kompositionen, Aufführungen und Veröffentlichungen wurde er zum Herold der Neuen Musik in Deutschland. Praetorius starb am 15. Februar 1621 und wurde unter der Orgelempore der Hauptkirche Beatae Mariae Virginis in Wolfenbüttel beigesetzt.
	
	\textit{Quelle: mh-koeln.de/cck}
	
	\section*{Randall Thompson}
	\paragraph{(geboren 1899 in New York City; gestorben 1984 in Boston) \newline}
	Randall Thompson war ein US-amerikanischer Komponist und Musikpädagoge, der vor allem für seine Chorwerke bekannt wurde. Er studierte an der Harvard University und am American Conservatory in Rom. 1927 wurde er Professor für Musik, Organist und Chorleiter am Wellesley College. Später unterrichtete er an der Harvard University, leitete verschiedene Madrigalchöre und den New Yorker Juilliard Choir und wurde Professor an der Berkeley University, der University of Virginia und der Princeton University. Seine Studie College Music (1935) bewirkte eine grundlegende Reform der universitären Musikausbildung in den USA. Thompson komponierte drei Sinfonien, zwei Opern, weitere sinfonische Werke und Kammermusik. Populär wurden jedoch vor allem seine Chorwerke.
	
	\textit{Quelle: Wikipedia}
	\vfill
	\pagebreak
	
	\section*{Camerata Vocale München}
	Die Camerata Vocale München gründete sich 2016 auf Initiative von Clayton Bowman zunächst als Projektchor mit Sängerinnen und Sängern aus München und Heidelberg.
	
	Als erstes Projekt gestaltete die Camerata am Palmsonntag 2016 einen Motettengottesdienstes mit Bachs „Jesu, meine Freude“  in der ältesten evangelischen Kirche Münchens, der St. Paulus-Kirche im Stadtteil Perlach. Im selben Jahr folgte ein Pfingstkonzert mit Chormusik aus fünf Jahrhunderten mit Werken von Reger, Purcell, Britten, Gesualdo, Whitacre, Vaughan Williams u.a., ebenfalls in St. Paulus/Perlach. Im März 2017 erarbeitete die Camerata Vocale mit dem Stabat Mater von Domenico Scarlatti, dem Requiem von Ildebrando Pizzetti und dem Crucifixus von Antonio Lotti ihr letztes Programm als Projektchor, das in der Nikodemuskirche München und in der Lutherkirche Mannheim anlässlich einer Glockeneinweihung aufgeführt wurde.
	
	Seit Sommer 2017 besteht die Camerata Vocale als ambitionierter Münchner Kammerchor mit regelmäßiger Probenarbeit und zählt inzwischen 23 Mitglieder. Nach kleineren Auftritten in Gottesdiensten der KHG der LMU und in St. Paulus/Perlach, gab die Camerata im Januar 2018 ebenda ihr Debütkonzert mit A-Capella-Werken u.a. von Reger, Brahms und Vaughan Williams. Im April 2018 wirkte sie an einer Produktion der Hochschule für Musik und Theater München mit, in deren Rahmen die Barockoper La Dafne von Marco da Gagliano im Rahmen der Barocktage 2018 im Großen Konzertsaal der HMT und in der Kreuzkirche Schwabing zur Aufführung gebracht wurde. Die Konzertsaison beschloss der Chor mit zwei Auftritten in St. Paulus in München/Perlach sowie St. Leonhard in Nußdorf am Inn im Juli 2018, bei denen jeweils anspruchsvolle Vokalmusik aus der englischen Chortradition mit Werken von Henry Purcell, Benjamin Britten, Edward Elgar, Ralph Vaughan Williams, Herbert Howells und William Walton sowie solistisch vorgetragene Lieder von Henry Purcell, John Dowland und Michael Tippett mit Continuo-Begleitung erklangen.
	
	Irgendein Text über die Adventskonzerte in Rosenheim und Perlach.
	
	\paragraph{Besetzung\newline}
	\begin{tabularx}\textwidth {lX}
	Sopran I & Susanne H. \newline Cosima Stocker \newline Sina Dresp \\
	Sopran II & Kathrin Sollfrank \newline Ruth Stärcke \\
	Alt I &  Anna Hausner \newline Esther Stärcke \newline Nora Reinbold \\
	Alt II & Salome Schuldt \newline Hannah S. \\
	Tenor I & Philipp Hummel \newline Julius Kiendl \\
	Tenor II & Benedikt Linder \newline Andru??? \\
	Bass I & Konrad Brückel \newline Tobias Gumpp \newline Freddy Zumegen \\
	Bass II & Joe Kübler \newline Andi Scharfstein
	\end{tabularx}
	\pagebreak

	\section*{Clayton Bowman}
	Clayton Bowman, geboren in Pittsburgh (USA), erhielt seinen ersten Klavierunterricht im Alter von 7 Jahren und wurde schon in jungen Jahren an die Chormusik herangeführt. Er sang in mehreren Auswahlchören und durfte bereits im Alter von 13 Jahren als Knaben-Solist sein Operndebüt in Benjamin Brittens "`Noye's Fludde"' mit dem Hartford Symphony Orchestra feiern.
	
	Erste Dirigiererfahrung sammelte Bowman während seiner Schulzeit. Bereits mit 14 Jahren leitete er als "`Student Conductor"' das symphonische Blasensemble an der Tolland Middle School und nahm im Jahr 2000 an einem Dirigiermeisterkurs an der University of South Carolina teil. Ab 2001 studierte er an der University of Connecticut Gesang und Musikwissenschaft mit Schwerpunkt Ensemblearbeit. Dieses Studium führte ihn schließlich nach Deutschland, wo er von 2004 bis 2008 Dirigieren bei Prof. Georg Grün, Klaus Thielitz und Wolfgang Seeliger an der Hochschule für Musik und Darstellende Kunst Mannheim studierte. Parallel dazu war Bowman lange Zeit Sänger und Dirigent zahlreicher Ensembles der Rhein-Neckar-Region.
	
	Als musikalischer Assistent beim Konzert Darmstadt und stellvertretender Dirigent des Universitätsorchesters Mannheim sammelte der junge Dirigent Erfahrung mit der Einstudierung großer Werke, wie beispielsweise Johann Sebastian Bachs "`Matthäuspassion"' und hatte zeitgleich die Leitung des Kammerchors Altrip, des Prot. Kirchenchors Mutterstadt und der Frauenchöre der Prot. Gemeinde in Dannstadt inne. Zusätzlich zu seinen Tätigkeiten als Dirigent, musizierte er häufig zusammen mit der Heidelberger Kantorei und der Hochschule für Kirchenmusik Heidelberg als Chorsänger und Solist. Darüber hinaus sang er regelmäßig im Kammerchor Saarbrücken, im Chor der Staatsphilharmonie Rheinland-Pfalz und gastierte gelegentlich als Chorist auf der Bühne des Nationaltheaters Mannheim und an der Oper Frankfurt.
	 
	Besonders aber mit dem Anglistenchor, einem der beiden Kammerchöre der Universität Heidelberg, hat sich Bowman einen Namen gemacht und eine große Leidenschaft für anspruchsvolle A-cappella-Musik entwickelt. Der Chor wirkte unter seiner Leitung vielfach im Ausland und arbeitete mit weltbekannten Chören wie dem Choir of Gonville und dem Caius College Cambridge zusammen.
	 
	Seit 2016 in München, leitet Bowman neben der vom ihm gegründeten Camerata Vocale München auch den Großen Chor des Akademischen Gesangvereins (AGV) und den Chor TonArt Sauerlach-Holzkirchen.
	\vfill
	\pagebreak
	
	\section*{Texte}
	
	\paragraph{Trois Chansons (1898 - 1908)\newline}
	I. Dieu qu'il la fait bon regarder!\newline
	La gracieuse bon et belle!\newline
	Pour les grands bien que sont en elle.\newline
	Chacun est pres de la loüer.\newline
	Qui se pourroit d'elle lasser?\newline
	Tousjours sa beauté renouvelle.\newline
	Dieu qu'il la bon regarder.\newline
	La gracieuse bonne et belle!\newline
	Par de ça, ne de là, la mer.\newline
	Ne scay dame ne damoiselle\newline
	Qui soit en tous bien parfais telle.\newline
	C'est un songe que d'i penser:\newline
	Dieu qu'il la fait bon regarder!\newline
	Dieu! qu'il la fait bon regarder!\newline
		
	II. Quant j'ai ouy le tabourin\newline
	Sonner pour s'en aller au may\newline
	En mon lit n'en ay fait effray\newline
	Ne levé mon chief du coissin\newline
	En disant: il est trop matin\newline
	Ung peu je me rendormirai:\newline
	Quant j'ai ouy le tabourin\newline
	Sonner pour s'en aller au may.\newline
	Jeunes gens partent leur butin:\newline
	De non chaloir m'accointeray\newline
	A lui je m'abutineray\newline
	Trouvé l'ay plus prouchain voisin;\newline
	Quant j'ai ouy le tabourin\newline
	Sonner pour s'en aller au may.\newline
	En mon lit n'en ay fait effray\newline
	Ne levé mon chief du coissin.\newline
	
	III. Yver, vous n'estes qu'un villain;\newline
	Esté est plaisant et gentil\newline
	En témoing de may et d'avril\newline
	Qui l'accompaignent soir et main.\newline
	Esté revet champs, bois et fleurs\newline
	De sa livrée de verdure\newline
	Et de maintes autres couleurs\newline
	Par l'ordonnance de nature.\newline
	Mais vous Yver, trop estes plein\newline
	De nège, vert, pluye et grézil.\newline
	On vous deust banir en évil.\newline
	Sans point flater je parle plein,\newline
	Yver, vous n'estes qu'un villain.\newline
	
	\paragraph{Hosanna to the Son of David\newline}
	Hosanna to the Son of David,\newline
	blessed be the King that cometh\newline
	in the name of the Lord, Hosanna.\newline
	Thou that sittest in the highest heav’ns,\newline
	Hosanna in excelsis Deo.
	
	\paragraph{Machet die Tore weit\newline}
	Machet die Tore weit\newline
	und die Türen in der Welt hoch,\newline
	daß der König der Ehren einziehe.\newline
	Wer ist der selbige König der Ehren?\newline
	Es ist der Herr,\newline
	stark und mächtig im Streit.\newline
	Machet die Tore weit\newline
	und die Türen in der Welt hoch!\newline
	Hosianna in der Höhe.\newline
	Hosianna dem Sohne Davids.
	
	\paragraph{In dulci jubilo\newline}
	In dulci jubilo,\newline
	Nun singet und seid froh!\newline
	Unsers Herzens Wonne\newline 
	Leit in praesepio\newline
	Und leuchtet als die Sonne,\newline
	Matris in gremio,\newline
	Alpha es et O.
	
	\paragraph{Alleluia\newline}
	Alleluia.
		
	\end{multicols}
	
	\pagebreak
	\centering
	\begin{figure}
		\includegraphics[width=\textwidth]{img/chorfoto_heller.png}
		\tiny{Foto: Chris Regel}
	\end{figure}
	\par\vspace{2cm}
	
	{\scshape\large Wir freuen uns sehr über Ihre Spende am Ende des Konzerts.\par
	Durch Ihre Unterstützung können wir unter anderem Programmhefte drucken oder Noten erwerben.\par}
	\vspace{1.5cm}
	{\scshape\Large Vielen Dank\par}
	\vfill
	{\small Impressum\par
		{\large Camerata Vocale München\par}
	Leitung: Clayton Bowman\par
	Kontakt: info@cameratavocale-muenchen.de\par
	www.cameratavocale-muenchen.de\par
	www.facebook.com/CVMuenchen\par
	Texte und Gestaltung: Cosima Stocker, Denise Dudek, Benedikt Linder}
	

\end{document}